\section{Технологическая часть}
\label{sec:tech}

\subsection{Инструкция пользователя}
\label{subsec:user-manual}

\subsubsection{Требования к среде}

Для сборки и запуска интерпретатора необходимо следующее программное обеспечение
(таблица~\ref{tab:requirements}). Проверка наличия и версий выполняется
командами \prog{java -\/\/-version}, \prog{mvn -\/\/-version} в терминале.

\begin{table}[!htb]
\caption{\centering Требования к среде разработки}
\label{tab:requirements}
\centering
\small
\begin{tabular}{@{}lll@{}}
\toprule
\textbf{Компонент} & \textbf{Минимальная версия} & \textbf{Примечание} \\
\midrule
Java Development Kit & 25 LTS & Сборка и запуск \\
Apache Maven & 3.9.x & Сборка проекта \\
Python~3 & 3.9+ & Только для построения графиков \\
\bottomrule
\end{tabular}
\end{table}

\subsubsection{Сборка}

Проект собирается командой Apache~Maven из корня репозитория \prog{refal5j-impl/}:

\begin{longlisting}
\caption{Команда сборки}\label{lst:build}
\begin{minted}{bash}
mvn package -DskipTests
\end{minted}
\end{longlisting}

Результат сборки --- файл \prog{target/refal5q.jar}, самодостаточный
JAR-архив с прописанным классом \prog{Main-Class: refal.Main}.
Все зависимости рантайма включены в него через плагин Maven~Shade,
поэтому для запуска достаточно стандартной JVM без дополнительных библиотек.

Сборка с полным прогоном тестов выполняется командой \prog{mvn package}
(без флага \prog{-DskipTests}) или \prog{mvn verify}, которая дополнительно
формирует отчёт о покрытии кода инструментом JaCoCo.

\newpage
\subsubsection{Запуск интерпретатора}

Интерпретатор запускается через стандартную команду \prog{java -jar}:

\begin{longlisting}
\caption{Форма запуска}\label{lst:run}
\begin{minted}{bash}
java -jar target/refal5q.jar [OPTIONS] program.ref
\end{minted}
\end{longlisting}

Позиционный аргумент --- путь к исходному файлу \prog{.ref}.
Поддерживаемые опции перечислены в таблице~\ref{tab:options}.

\begin{table}[!htb]
\caption{\centering Опции командной строки}
\label{tab:options}
\centering
\small
\begin{tabular}{@{}lp{8.5cm}@{}}
\toprule
\textbf{Опция} & \textbf{Описание} \\
\midrule
\prog{--entry <имя>} &
    Явное указание входной функции; имеет приоритет над \prog{\$ENTRY} в файле \\
\prog{--step-limit <N>} &
    Максимальное число шагов нормализации (по умолчанию~--- 100\,000) \\
\prog{--expr <терм...>} &
    Начальное выражение, передаваемое входной функции вместо пустого \\
\bottomrule
\end{tabular}
\end{table}

Код завершения процесса отражает результат выполнения
(таблица~\ref{tab:exit-codes}).

\begin{table}[!htb]
\caption{\centering Коды завершения процесса}
\label{tab:exit-codes}
\centering
\small
\begin{tabular}{@{}cl@{}}
\toprule
\textbf{Код} & \textbf{Условие} \\
\midrule
0 & Успешное завершение или \prog{<Exit 0>} \\
1 & Ошибка лексического, синтаксического или семантического анализа \\
2 & Ошибка времени выполнения: нет подходящего правила,
    превышение лимита шагов \\
$n$ & Значение $n$, переданное встроенной функции \prog{<Exit~n>} \\
\bottomrule
\end{tabular}
\end{table}

Пример запуска программы \prog{hello.ref} с явным указанием точки входа:

\begin{longlisting}
\caption{Пример запуска}\label{lst:run-example}
\begin{minted}{bash}
java -jar target/refal5q.jar --entry Go hello.ref
\end{minted}
\end{longlisting}

\subsubsection{Запуск тестов}

Полный набор тестов запускается одной командой Maven:

\newpage
\begin{longlisting}
\caption{Запуск тестов}\label{lst:test}
\begin{minted}{bash}
mvn test
mvn verify
mvn test -Dtest=AutotestsRunnerTest
\end{minted}
\end{longlisting}

Первые две команды прогоняют 143~JUnit-теста, покрывающих все компоненты
проекта. Третья команда запускает отдельный класс \prog{AutotestsRunnerTest},
который вызывает интерпретатор на 81~внешней программе из двух наборов:
автотесты репозитория \prog{refal-5j} и десахаренные тесты репозитория
\prog{refal-5-framework}. Каждая программа выполняется в отдельном процессе
с тайм-аутом 20~секунд.

Отчёт JaCoCo формируется в каталоге \prog{target/site/jacoco/}; открыть
его можно браузером, указав путь к \prog{index.html}.

%%% ─────────────────────────────────────────────────────────
\subsection{Структура данных: раскрученная двусторонняя очередь}
\label{subsec:deque-impl}

Объектное выражение в нашей реализации хранится в раскрученной (bootstrapped)
двусторонней очереди по Окасаки~\cite{okasaki_pfds}. У неё амортизированно
константная конкатенация~--- та самая операция, на которой классические
списочные интерпретаторы Рефала (включая Refal-5~PZ и Refal-5J) тратят
линейное время.

\subsubsection{Sealed interface Expr<T>}

Контейнер выражения объявлен как обобщённый запечатанный интерфейс с двумя
реализациями:

\begin{longlisting}
\caption{Запечатанный интерфейс Expr<T>}\label{lst:expr-iface}
\begin{minted}{java}
public sealed interface Expr<T> extends Iterable<T>
        permits ShallowDeque, DeepDeque {
    boolean isEmpty();
    int size();
    PeekResult<T> separateFront();
    PeekResult<T> separateBack();
    Expr<T> pushFront(T x);
    Expr<T> pushBack(T x);
    static <T> Expr<T> empty()        { return ShallowDeque.empty(); }
    static <T> Expr<T> singleton(T x) { return ShallowDeque.<T>empty().pushBack(x); }
    static <T> Expr<T> concat(Expr<T> a, Expr<T> b) { return ExprOps.concat(a, b); }
}
\end{minted}
\end{longlisting}

Слово \prog{sealed} запрещает любые другие наследники. Компилятор
Java~25 на этом основании проверяет полноту разбора в
\prog{switch}-выражениях, поэтому в коде конкатенации мы можем
перечислить ровно четыре случая по типам операндов и не писать ветку
\prog{default}: её отсутствие гарантировано системой типов.

\subsubsection{Параметры буферов}

Ёмкость мелкого буфера~--- константа \prog{ShallowDeque.CAPACITY = 9},
нижняя граница заполненности \prog{MIN\_FILL = 5}, то есть
$\lceil CAPACITY / 2 \rceil$. Нечётное значение выбрано по
арифметическим соображениям: при переполнении мелкого буфера суммарно
получается $CAPACITY + 1 = 10$ элементов; они делятся ровно пополам
на два фрагмента по 5, и оба фрагмента сразу удовлетворяют
\prog{MIN\_FILL}. С чётной ёмкостью пришлось бы либо мириться с
остатком в один элемент, либо добавлять отдельную обработку этого
случая~--- ни то ни другое не даёт выигрыша по свойствам структуры.

На практике глубина рекурсии \prog{DeepDeque} при таком \prog{CAPACITY}
почти не выходит за два уровня: в тестовых программах с выражениями
длиной до десятков тысяч термов средняя очередь сама ни разу не успевает
стать глубокой.

\subsubsection{ShallowDeque<T>: мелкий буфер}

\prog{ShallowDeque<T>}~--- иммутабельный класс с тремя полями:
массив фиксированной длины \prog{Object[] buf} размером \prog{CAPACITY},
индекс головы \prog{head} и текущий размер \prog{sz}.
Доступ по логическому индексу~$i$ выполняется через единственную формулу:

\newpage
\begin{longlisting}
\caption{Циклическое индексирование в ShallowDeque}\label{lst:shallow-idx}
\begin{minted}{java}
int idx(int i) {
    return Math.floorMod(head + i, CAPACITY);
}
\end{minted}
\end{longlisting}

\prog{floorMod} здесь стоит вместо обычного \prog{\%}, потому что в Java
оператор \prog{\%} при отрицательном делимом возвращает отрицательный
остаток, а нам нужен индекс в массиве. \prog{pushFrontTight} как раз
сдвигает голову на~$-1$, и без \prog{floorMod} в этой точке пришлось бы
вручную добавлять \prog{CAPACITY}.

В классе определены два варианта операций добавления:
\begin{itemize}
  \item \prog{pushFrontTight(T)} и \prog{pushBackTight(T)}~--- внутренние, возвращают
        \prog{ShallowDeque<T>} и бросают \prog{IllegalStateException} при переполнении;
        применяются там, где вызывающий код гарантирует наличие свободного места
        (например, при формировании буферов во время конкатенации);
  \item \prog{pushFront(T)} и \prog{pushBack(T)}~--- публичные (методы
        интерфейса \prog{Expr<T>}); при переполнении автоматически переходят
        в глубокий уровень через \prog{DeepDeque.mkDeep}.
\end{itemize}

\subsubsection{DeepDeque<T>: глубокий уровень}

\prog{DeepDeque<T>} реализует тройку $\langle front, mid, back \rangle$:

\begin{longlisting}
\caption{Поля DeepDeque<T>}\label{lst:deep-fields}
\begin{minted}{java}
public final class DeepDeque<T> implements Expr<T> {
    private final ShallowDeque<T>          front;
    private final Expr<ShallowDeque<T>>    mid;
    private final ShallowDeque<T>          back;
    private final int                      sz;
    ...
}
\end{minted}
\end{longlisting}

Тип поля \prog{mid}~--- \prog{Expr<ShallowDeque<T>>}, и именно здесь
проявляется «раскрутка»: средняя очередь хранит мелкие буферы, но сама
она тоже \prog{Expr<\ldots>}, то есть может оказаться как
\prog{ShallowDeque}, так и \prog{DeepDeque} следующего уровня.
На разумных размерах выражений рекурсия редко уходит глубже второго
уровня.

Конструктор \prog{mkDeep} автоматически схлопывает глубокую очередь обратно
в мелкую, когда суммарное число элементов не превышает \prog{CAPACITY}:

\begin{longlisting}
\caption{Конструктор mkDeep с откатом в ShallowDeque}\label{lst:mkdeep}
\begin{minted}{java}
static <T> Expr<T> mkDeep(ShallowDeque<T> front,
                           Expr<ShallowDeque<T>> mid,
                           ShallowDeque<T> back, int total) {
    if (total <= CAPACITY) {
        ShallowDeque<T> flat = ShallowDeque.empty();
        for (T x : iterableOf(front))  flat = flat.pushBackTight(x);
        for (ShallowDeque<T> bucket : iterableOf(mid))
            for (T x : iterableOf(bucket)) flat = flat.pushBackTight(x);
        for (T x : iterableOf(back))   flat = flat.pushBackTight(x);
        return flat;
    }
    return new DeepDeque<>(front, mid, back, total);
}
\end{minted}
\end{longlisting}

Параметр \prog{total} передаётся явно, чтобы не пересчитывать его внутри:
\prog{mid.size()} возвращает число буферов в средней очереди, а не суммарное
число листовых элементов, поэтому наивная формула
\prog{front.size() + mid.size() + back.size()} оказалась бы некорректной.
Каждый вызов \prog{mkDeep} обновляет \prog{total} инкрементально~---
прибавляя или вычитая единицу по ходу операции добавления или извлечения.

\subsubsection{Конкатенация: четыре случая}

Конкатенация реализована в служебном классе \prog{ExprOps} с разбором
четырёх случаев по мелкости/глубине операндов:

\begin{longlisting}
\caption{Разбор четырёх случаев конкатенации}\label{lst:concat-cases}
\begin{minted}{java}
static <T> Expr<T> concat(Expr<T> a, Expr<T> b) {
    if (a.isEmpty()) return b;
    if (b.isEmpty()) return a;
    if (a instanceof ShallowDeque<T> sa && b instanceof ShallowDeque<T> sb)
        return concatSS(sa, sb);
    if (a instanceof ShallowDeque<T> sa && b instanceof DeepDeque<T> db)
        return concatSD(sa, db);
    if (a instanceof DeepDeque<T> da && b instanceof ShallowDeque<T> sb)
        return concatDS(da, sb);
    DeepDeque<T> da = (DeepDeque<T>) a;
    DeepDeque<T> db = (DeepDeque<T>) b;
    return concatDD(da, db);
}
\end{minted}
\end{longlisting}

В случае \emph{Deep\,+\,Deep} пограничные буферы \prog{a.back} и
\prog{b.front} склеиваются в один буфер, а \prog{concat} вызывается
рекурсивно уже на средних очередях. На каждом уровне рекурсии работа
ограничена константой, а сама рекурсия идёт не глубже логарифма от
длины выражения~--- так Окасаки и получает амортизированную
константу~\cite{okasaki_pfds}.

%%% ─────────────────────────────────────────────────────────
\subsection{Цикл нормализации}
\label{subsec:runtime-impl}

\subsubsection{Поле зрения и иерархия ViewTerm}

Нормализация выражения выполняется итеративным проходом по
\emph{полю зрения}~--- линейной последовательности элементарных единиц.
Каждая такая единица описывается запечатанным интерфейсом \prog{ViewTerm}:

\begin{longlisting}
\caption{Иерархия ViewTerm}\label{lst:viewterm}
\begin{minted}{java}
public sealed interface ViewTerm
        permits ViewTerm.Passive,  ViewTerm.BulkPassive,
                ViewTerm.AngOpen,  ViewTerm.AngClose,
                ViewTerm.ParenOpen, ViewTerm.ParenClose {

    record Passive(Term term)          implements ViewTerm {}
    record BulkPassive(Expr<Term> items) implements ViewTerm {}
    record AngOpen(String fn)          implements ViewTerm {}
    record AngClose()                  implements ViewTerm {}
    record ParenOpen()                 implements ViewTerm {}
    record ParenClose()                implements ViewTerm {}
}
\end{minted}
\end{longlisting}

Термы-атомы (\prog{Sym}, \prog{Num}, \prog{Str}) и скобочные термы \prog{Br}
после нормализации вложенных выражений представляются как \prog{Passive}.
Вызов \prog{<f args>} разворачивается в последовательность
\prog{AngOpen(f)}, элементы аргументов, \prog{AngClose}.
Скобочный терм~\prog{(e.X)} аналогично разворачивается через
\prog{ParenOpen} и \prog{ParenClose}.
\prog{BulkPassive} отвечает за оптимизацию связывания переменных и описан
в следующем подразделе.

\subsubsection{Двустековая реализация нормализатора}

Нормализация работает с двумя стеками~--- \prog{left} и \prog{right}:
\begin{itemize}
  \item \prog{right} содержит ещё не обработанные элементы поля зрения;
  \item \prog{left} накапливает уже пассивные элементы~--- те, в которых
        не осталось ни одного активного вызова.
\end{itemize}

Исходное выражение сначала сплющивается в \prog{right} слева-направо.
Затем главный цикл извлекает элементы из \prog{right} по одному с головы.
Пассивные элементы (\prog{Passive}, \prog{BulkPassive}) и открывающие
маркеры (\prog{AngOpen}, \prog{ParenOpen}) просто перекладываются в
\prog{left}. Обработка наступает при встрече закрывающего маркера:

\begin{longlisting}
\caption{Обработка закрывающего маркера вызова функции}\label{lst:angclose}
\begin{minted}{java}
case ViewTerm.AngClose ignored -> {
    bumpStep();
    Expr<Term> argsExpr = Expr.empty();
    String fn = null;
    while (true) {
        ViewTerm top = left.pollLast();
        if (top instanceof ViewTerm.AngOpen ao) { fn = ao.fn(); break; }
        if (top instanceof ViewTerm.Passive pt)
            argsExpr = argsExpr.pushFront(pt.term());
        else if (top instanceof ViewTerm.BulkPassive bp)
            argsExpr = Expr.concat(bp.items(), argsExpr); // O(1)
    }
}
\end{minted}
\end{longlisting}

При встрече \prog{AngClose} мы снимаем элементы с хвоста \prog{left}
до парного \prog{AngOpen} и попутно собираем выражение аргументов
\prog{argsExpr}. Одиночный \prog{Passive} присоединяется через
\prog{pushFront} за~$O(1)$, а \prog{BulkPassive}~--- через
\prog{Expr.concat(bp.items(), argsExpr)}, тоже за~$O(1)$
амортизированно. Обработка \prog{ParenClose} устроена точно так же,
только без поиска имени функции.

После сборки аргументов для встроенной функции результат сплющивается и
помещается обратно в \prog{right}; для пользовательской функции находится
первое подходящее правило и его правая часть, с подставленными связываниями
переменных, укладывается на \prog{right} методом
\prog{pushResultToStack}.

Итеративный цикл вместо рекурсивного обхода даёт нам независимость от
размера стека JVM: сколь угодно глубоко вложенная рефал-программа
обрабатывается в константной глубине Java-стека, а реальный расход
памяти зависит только от длины поля зрения.

\subsubsection{Оптимизация BulkPassive}
\label{subsubsec:bulkpassive}

Без оптимизации подстановка значения переменной в правой части правила
требовала поэлементной укладки: каждый \prog{Term} из связанного значения
помещался на стек \prog{right} как отдельный объект \prog{Passive}:

\begin{longlisting}
\caption{Наивная укладка переменной (до оптимизации)}\label{lst:var-naive}
\begin{minted}{java}
case refal.ast.Var v -> {
    Expr<Term> val = bindings.get(v.key());
    // O(N) итераций: каждый терм становится Passive
    for (Term t : val) {
        right.addFirst(new ViewTerm.Passive(t));
    }
}
\end{minted}
\end{longlisting}

При следующей встрече \prog{AngClose} те же $N$ элементов снимались
с \prog{left} по одному~--- ещё $O(N)$ работы на каждый шаг нормализации.
В программе, которая на каждом шаге накапливает выражение длиной $N$
(как Loop/Select из аналитической части), это давало $O(N^2)$
операций суммарно.

Оптимизация сводится к одному изменению: вместо перебора элементов
переменной её значение помещается на стек одним объектом \prog{BulkPassive}:

\begin{longlisting}
\caption{Укладка переменной через BulkPassive (O(1))}\label{lst:var-bulk}
\begin{minted}{java}
case refal.ast.Var v -> {
    Expr<Term> val = bindings.get(v.key());
    if (!val.isEmpty()) {
        right.addFirst(new ViewTerm.BulkPassive(val)); // O(1)
    }
}
\end{minted}
\end{longlisting}

Когда обработчик \prog{AngClose} снимает этот объект с \prog{left},
он подмешивает его в накапливаемые аргументы вызовом \prog{Expr.concat}~---
тоже амортизированно $O(1)$ (листинг~\refAlgo{lst:angclose}).

Важный инвариант: \prog{BulkPassive} попадает в поле зрения ровно в той
позиции, где раньше лежали бы отдельные \prog{Passive} с тем же набором
термов в том же порядке. Это значит, что замена не меняет порядок
вычисления и не нарушает семантику Рефала~--- мы просто переносим
ту же подпоследовательность одним объектом вместо $N$.

Эффект от оптимизации хорошо виден на двух тестах. Тест
\prog{03-concat.ref}, который раньше упирался в лимит в 100~миллионов
шагов нормализации и не завершался, теперь проходит примерно за
10~секунд. На бенчмарке Loop/Select при $N = 20\,000$ наш интерпретатор
показывает порядка 300~мс против секунды и более у Refal-5~PZ и
Refal-5J.
